\section*{第 \ref{ch:g8:pythagoras} 章 --- 勾股定理}

\begin{solution}{exo:g8:pythagoras:1}
$\sqrt{81 + 144} = \sqrt{225} = 15$; \quad
$\sqrt{25 + 144} = \sqrt{169} = 13$; \quad
$\sqrt{64 + 225} = \sqrt{289} = 17$.
\end{solution}

\begin{solution}{exo:g8:pythagoras:2}
直角边 $^2 = 25^2 - 7^2 = 625 - 49 = 576$，故该直角边为
$\sqrt{576} = 24$。
\end{solution}

\begin{solution}{exo:g8:pythagoras:3}
长方形：$\sqrt{12^2 + 9^2} = \sqrt{144 + 81} = \sqrt{225} = 15$
cm。

正方形：$d = 5\sqrt2 \approx 7.1$ cm
（见\cref{ex:g8:pythagoras:diagonal}）。
\end{solution}

\begin{solution}{exo:g8:pythagoras:4}
$(20, 21, 29)$：$29^2 = 841$ 且 $20^2 + 21^2 = 400 + 441 = 841$：
为直角三角形。

$(7, 8, 11)$：$121 \neq 49 + 64 = 113$：不是直角三角形。

$(1.5, 2, 2.5)$：$6.25 = 2.25 + 4$：为直角三角形（它是
$(3,4,5)$ 各边减半所得）。
\end{solution}

\begin{solution}{exo:g8:pythagoras:5}
木板长 $= \sqrt{3^2 + 1.6^2} = \sqrt{9 + 2.56} = \sqrt{11.56} = 3.4$
m。
\end{solution}

\begin{solution}{exo:g8:pythagoras:6}
高落在底边中点上（等腰三角形），故
每个直角三角形的斜边为 $10$，一条直角边为 $6$：
$h = \sqrt{100 - 36} = \sqrt{64} = 8$ cm。面积：
$\frac{12 \times 8}{2} = 48$ cm$^2$。
\end{solution}

\begin{solution}{exo:g8:pythagoras:7}
对角线：$\sqrt{88.5^2 + 49.8^2} = \sqrt{7832.25 + 2480.04}
= \sqrt{10312.29} \approx 101.5$ cm。换算为英寸：
$101.5 \div 2.54 \approx 40$ 英寸。
\end{solution}

\begin{solution}{exo:g8:pythagoras:8}
两直角边分别为 $24$ km 和 $10$ km：
$\sqrt{24^2 + 10^2} = \sqrt{576 + 100} = \sqrt{676} = 26$ km。
\end{solution}

\begin{solution}{exo:g8:pythagoras:9}
梯脚到墙的距离：$\sqrt{2.5^2 - 2.4^2} = \sqrt{6.25 - 5.76} =
\sqrt{0.49} = 0.7$ m。梯子长度的四分之一为
$2.5 \div 4 \approx 0.6$ m：求得的 $0.7$ m 接近建议的
放置位置 --- 可以接受。
\end{solution}

\begin{solution}{exo:g8:pythagoras:10}
两直角边长分别为 $5 - 1 = 4$（水平）和 $5 - 2 = 3$
（竖直方向），故
\[
AB = \sqrt{4^2 + 3^2} = \sqrt{25} = 5 .
\]
\end{solution}

\begin{solution}{exo:g8:pythagoras:11}
\emph{1.} 直接看：$(a + b)^2$。拆开来看：四个面积各为
$\frac{ab}{2}$ 的三角形，再加上倾斜放置的正方形 $c^2$：
\[
(a + b)^2 = 4 \times \frac{ab}{2} + c^2 = 2ab + c^2 .
\]

\emph{2.} 展开左边：
$(a+b)^2 = a^2 + 2ab + b^2$。于是
\[
a^2 + 2ab + b^2 = 2ab + c^2
\quad\Longrightarrow\quad
a^2 + b^2 = c^2 . \qed
\]
\end{solution}

\begin{solution}{pb:g8:pythagoras:1}
\textbf{1.} $5^2 + 12^2 = 25 + 144 = 169 = 13^2$; \quad
$8^2 + 15^2 = 64 + 225 = 289 = 17^2$; \quad
$20^2 + 21^2 = 400 + 441 = 841 = 29^2$。三组数均为勾股数，
由\cref{thm:g8:pythagoras:converse}可知，三边长为这些数的三角形
为直角三角形，其直角位于最长
边所对的角。

\textbf{2.} 绳子围成一个三角形，三边分别为 $3$、$4$、$5$，
（以段计），且 $3^2 + 4^2 = 25 = 5^2$：由逆定理
（\cref{thm:g8:pythagoras:converse}），两边分别为 $3$ 段与 $4$ 段，其夹角
恰为直角 --- 无需量角器，
只需要打结。

\textbf{3.} 若 $a^2 + b^2 = c^2$，则
\[
(ka)^2 + (kb)^2 = k^2 a^2 + k^2 b^2
= k^2\left(a^2 + b^2\right) = k^2 c^2 = (kc)^2 .
\]
由 $(3, 4, 5)$ 出发，例如 $(6, 8, 10)$、$(9, 12, 15)$ 和
$(30, 40, 50)$。

\textbf{4.} 若 $(3k, 4k, 5k)$ 的某个倍数以 $5$ 为其最小数，
则需 $3k = 5$，但没有整数 $k$
满足该式（$k = 1$ 时为 $3$，$k = 2$ 时为 $6$）。因此
$(5, 12, 13)$ 不是 $(3, 4, 5)$ 的倍数。

\textbf{5.} 第 3 题表明，放缩只能产生
形如 $(3k, 4k, 5k)$ 的三元组；第 4 题给出一个不属于此形式的
真正的勾股数。因此，全体勾股数组成的族
严格大于 $(3, 4, 5)$ 的倍数组成的族：要得到它们
需要新的方法。

\textbf{6.} 记 $A = m^2$，$B = n^2$。由乘法公式
（\cref{pb:g8:equations:1}）可得：
\[
a^2 = (A - B)^2 = A^2 - 2AB + B^2,
\qquad
b^2 = (2mn)^2 = 4m^2n^2 = 4AB .
\]
两式相加：
\[
a^2 + b^2 = A^2 - 2AB + B^2 + 4AB
= A^2 + 2AB + B^2 = (A + B)^2
= \left(m^2 + n^2\right)^2 = c^2 .
\]

\textbf{7.}
\begin{center}
\renewcommand{\arraystretch}{1.2}
\begin{tabular}{c|ccc}
$(m, n)$ & $m^2 - n^2$ & $2mn$ & $m^2 + n^2$ \\
\hline
$(2, 1)$ & $3$ & $4$ & $5$ \\
$(3, 1)$ & $8$ & $6$ & $10$ \\
$(3, 2)$ & $5$ & $12$ & $13$ \\
$(4, 1)$ & $15$ & $8$ & $17$ \\
$(4, 3)$ & $7$ & $24$ & $25$ \\
\end{tabular}
\end{center}

\textbf{8.} $(8, 6, 10)$ 是 $(4, 3, 5)$ 的两倍 --- 也就是
建造者的 $(3, 4, 5)$。机器有时会以伪装的形式
重新找到一组旧的勾股数。

\textbf{9.} 需要 $m^2 - n^2 = 21$ 且 $2mn = 20$，因此
$mn = 10$：取 $m = 5$，$n = 2$，得 $25 - 4 = 21$ 且
$2 \times 10 = 20$，斜边为 $25 + 4 = 29$。机器
输出 $(21, 20, 29)$，对应 $(m, n) = (5, 2)$。

\textbf{10.} 取 $m = 8$，$n = 6$，得 $m^2 + n^2 = 64 + 36 = 100$。
所得勾股数为
\[
\left(64 - 36,\ 2 \times 48,\ 100\right) = (28,\ 96,\ 100),
\]
且确实 $28^2 + 96^2 = 784 + 9216 = 10\,000 = 100^2$。

\textbf{11.} 取 $m = n + 1$：
\[
m^2 - n^2 = (n + 1)^2 - n^2 = 2n + 1
\]
（相邻平方数之差，见\cref{pb:g8:equations:1}）；
$2mn = 2n(n + 1) = 2n^2 + 2n$；又
$m^2 + n^2 = n^2 + n^2 + 2n + 1 = 2n^2 + 2n + 1$。斜边
比偶直角边恰多 $1$。当 $n = 1, 2, 3, 4$ 时：
\[
(3, 4, 5), \quad (5, 12, 13), \quad (7, 24, 25), \quad
(9, 40, 41) .
\]

\textbf{12.} 每个不小于 $3$ 的奇数都是 $2n + 1$，其中
$n \geq 1$，而第 11 题给出一个三元组，其奇直角边恰为
$2n + 1$。对于 $11 = 2 \times 5 + 1$，取 $n = 5$：
\[
(11,\ 60,\ 61),
\qquad
11^2 + 60^2 = 121 + 3600 = 3721 = 61^2 .
\]

\textbf{13.} 取 $n = 1$，机器给出
$\left(m^2 - 1,\ 2m,\ m^2 + 1\right)$，其偶直角边 $2m$ 可取
从 $4$ 起的任一偶数。对于 $14 = 2 \times 7$，取
$m = 7$：得 $(48, 14, 50)$。同一组勾股数也可由
第 7 题表中的 $(7, 24, 25)$ 加倍得到 $(14, 48, 50)$
--- 即同样的三个数。

\textbf{14.} 偶数可写为 $2k$，而 $(2k)^2 = 4k^2$ 是偶数
（$4$ 的倍数，偶数）。奇数可写为 $2k + 1$，于是
由第一个乘法公式可得
\[
(2k + 1)^2 = 4k^2 + 4k + 1 = 2\left(2k^2 + 2k\right) + 1 ,
\]
即偶数加 $1$，是奇数。

\textbf{15.} 假设 $a$、$b$、$c$ 全为奇数。由第 14 题
可知 $a^2$ 与 $b^2$ 均为奇数，因此 $a^2 + b^2$ 是两个奇数之和：
\emph{偶数}。但 $c^2$ 应为 \emph{奇数}。等式
$a^2 + b^2 = c^2$ 于是不可能成立。因此在每个勾股数中，
三个数中至少有一个是偶数 --- 正如第 7 题表格
的每一行所印证的。
\end{solution}
